\documentclass[10pt,a4paper]{amsart}
\usepackage[margin=19mm]{geometry}
\usepackage{amsmath,amssymb,mathtools}
\usepackage[scr=rsfs,cal=euler]{mathalfa}
\usepackage{xfrac}
\usepackage[unicode,hidelinks,bookmarks=false]{hyperref}
\newcommand{\ie}{i.e.\ }
\newcommand{\cf}{cf.\ }
\newcommand{\resp}{respectively\ }
\newcommand{\Subsection}{\subsection}
\providecommand{\nobreakdash}{\nobreak\hbox{-}}
\setlength{\emergencystretch}{2em}
\multlinegap0pt
\usepackage{tikz}
\usepackage{amssymb}
\usepackage{enumerate}
\usepackage[shortlabels]{enumitem}
\newtheorem{proposition}{Proposition}
\newtheorem{lemma}{Lemma}
\def\XXint#1#2#3{{\setbox0=\hbox{$#1{#2#3}{\int}$ }
\vcenter{\hbox{$#2#3$ }}\kern-.6\wd0}}
\def\Xint#1{\mathchoice
{\XXint\displaystyle\textstyle{#1}}%
{\XXint\textstyle\scriptstyle{#1}}%
{\XXint\scriptstyle\scriptscriptstyle{#1}}%
{\XXint\scriptscriptstyle\scriptscriptstyle{#1}}%
\!\int}
\def\dashint{\Xint-}
\title{{On singularity of $p$-energy measures on metric measure spaces}}
\author{Meng Yang}
\date{Source: arXiv:2505.12468v2 (CC BY 4.0)}
\begin{document}
\maketitle

\noindent\emph{Source and licence.} {On singularity of $p$-energy measures on metric measure spaces}. Version arXiv:2505.12468v2 (CC BY 4.0), distributed by arXiv under the Creative Commons Attribution 4.0 International licence, http://creativecommons.org/licenses/by/4.0/, as recorded in the arXiv licence metadata for this exact version and shown on the abstract page https://arxiv.org/abs/2505.12468v2. Full licence notice: \texttt{licenses/arxiv-2505.12468-CC-BY-4.0.txt}.\par\smallskip

\noindent\textit{Editorial scope.} Counted unit: the article's proposition prop\_CS (source0.tex lines 601--603). The article's complete original proof of it is delivered with it (source0.tex lines 650--673, one \texttt{proof} environment of 24 lines, which the article places away from the statement). No mathematics is written, repaired or re-derived: the statement and the proof are the article's own lines, in the article's own notation, and no retained line is substituted or re-flowed.  Retained uncounted as the source-local context the proof needs: lines 140--143; lines 220--222; lines 228--230; lines 259--261; lines 607--611.  Nothing was omitted from any retained span, so the package carries no omission record.  No retained line contains a citation, so the wrapper carries no bibliography and no bibliography entry.  No retained line contains a figure environment, an image-inclusion command or drawing code, so no image is delivered and the package declares no asset dependency.  Editorial formatting only, in addition: an AMS \texttt{amsart} preamble that restates the article's own declarations and macros used by the retained lines, namely the environments \texttt{proposition}, \texttt{lemma} on separate counters, together with the standard packages those lines need.  The retained environments print in this document as Proposition 1, Lemma 1 in order of appearance, whereas the article prints the same items with the numbering of its own full document.  The complete native source of the article as distributed in the arXiv e-print for this exact version is bundled unchanged as \texttt{sources/178311/source0.tex}.
\begin{equation*}\label{eq_SVR}\tag*{$\text{SVR}(\Phi,\Psi)$}
\frac{\Phi(R)}{\Phi(r)}\le\Theta\left(\frac{r}{R}\right)\frac{\Psi(R)}{\Psi(r)}
\text{ for any }R,r\in(0,\mathrm{diam}(X))\text{ with }r\le R.
\end{equation*}

\begin{equation*}\label{eq_PI}\tag*{$\text{PI}(\Psi)$}
\int_B\lvert f-f_B\rvert^p\mathrm{d} m\le C_{PI}\Psi(r)\int_{A_{PI}B}\mathrm{d}\Gamma(f).
\end{equation*}

\begin{equation*}\label{eq_CS}\tag*{$\text{CS}(\Psi)$}
\int_{B(x,A_{S}r)}|\widetilde{f}|^p \mathrm{d}\Gamma(\phi)\le C_{1}\int_{B(x,A_{S}r)}\mathrm{d}\Gamma(f)+\frac{C_{2}}{\Psi(r)}\int_{B(x,A_{S}r)}|f|^p\mathrm{d} m,
\end{equation*}

\begin{equation*}\label{eq_cap}\tag*{$\text{cap}(\Psi)_{\le}$}
\mathrm{cap}\left(B(x,r),X\backslash B(x,A_{cap}r)\right)\le C_{cap} \frac{V(x,r)}{\Psi(r)}.
\end{equation*}

\begin{lemma}[Morrey-Sobolev inequality]\label{lem_MS}
Assume \ref{eq_SVR}, \ref{eq_PI}. Then there exists $C_{MS}>0$ such that for any $u\in\mathcal{F}$, for any ball $B(x_0,R)$, for $m$-a.e. $x,y\in B(x_0,R)$ with $x\ne y$, we have
$$|u(x)-u(y)|^p\le C_{MS} \frac{\Psi(d(x,y))}{\Phi(d(x,y))}\int_{B(x_0,8A_{PI}R)}\mathrm{d}\Gamma(u),$$
where $A_{PI}$ is the constant in \ref{eq_PI}. Hence any function in $\mathcal{F}$ has a continuous version, or equivalently, $\mathcal{F}\subseteq C(X)$.
\end{lemma}

\begin{proposition}\label{prop_CS}
Assume \ref{eq_SVR}, \ref{eq_PI}, \ref{eq_cap}. Then \ref{eq_CS} holds.
\end{proposition}

\begin{proof}[Proof of Proposition \ref{prop_CS}]
For any ball $B(x_0,R)$, by \ref{eq_cap}, there exists a cutoff function $\phi\in\mathcal{F}$ for $B(x_0,R)\subseteq B(x_0,A_{cap}R)$ such that
$$\mathcal{E}(\phi)=\int_{B(x_0,A_{cap}R)}\mathrm{d}\Gamma(\phi)\le2C_{cap}\frac{V(x_0,R)}{\Psi(R)}.$$
By Lemma \ref{lem_MS}, for any $f\in\mathcal{F}\subseteq C(X)$, for any $x,y\in B(x_0,A_{cap}R)$, by \ref{eq_SVR}, we have
\begin{align*}
&|f(x)-f(y)|^p\le C_{MS}\frac{\Psi(d(x,y))}{\Phi(d(x,y))}\int_{B(x_0,8A_{PI}A_{cap}R)}\mathrm{d}\Gamma(f)\\
&\le C_{MS}\Theta(1)\frac{\Psi(2A_{cap}R)}{\Phi(2A_{cap}R)}\int_{B(x_0,8A_{PI}A_{cap}R)}\mathrm{d}\Gamma(f),
\end{align*}
hence
\begin{align*}
&|f(x)|^p\le2^{p-1}(|f(x)-f(y)|^p+|f(y)|^p)\\
&\le2^{p-1}\left(C_{MS}\Theta(1)\frac{\Psi(2A_{cap}R)}{\Phi(2A_{cap}R)}\int_{B(x_0,8A_{PI}A_{cap}R)}\mathrm{d}\Gamma(f)+|f(y)|^p\right).
\end{align*}
Integrating both sides over $B(x_0,A_{cap}R)$ in the variable $y$ and dividing by $V(x_0,A_{cap}R)$, we obtain that for any $x\in B(x_0,A_{cap}R)$,
$$|f(x)|^p\le 2^{p-1}\left(C_{MS}\Theta(1)\frac{\Psi(2A_{cap}R)}{\Phi(2A_{cap}R)}\int_{B(x_0,8A_{PI}A_{cap}R)}\mathrm{d}\Gamma(f)+\dashint_{B(x_0,A_{cap}R)}|f|^p\mathrm{d} m\right).$$
Therefore, $\phi\in\mathcal{F}$ is also a cutoff function for $B(x_0,R)\subseteq B(x_0,8A_{PI}A_{cap}R)$, and by the strongly local property, we have
\begin{align*}
&\int_{B(x_0,8A_{PI}A_{cap}R)}|f|^p\mathrm{d}\Gamma(\phi)=\int_{B(x_0,A_{cap}R)}|f|^p\mathrm{d}\Gamma(\phi)\\
&\le2^{p-1}\left(C_{MS}\Theta(1)\frac{\Psi(2A_{cap}R)}{\Phi(2A_{cap}R)}\int_{B(x_0,8A_{PI}A_{cap}R)}\mathrm{d}\Gamma(f)+\dashint_{B(x_0,A_{cap}R)}|f|^p\mathrm{d} m\right)\\
&\hspace{20pt}\cdot\int_{B(x_0,A_{cap}R)}\mathrm{d}\Gamma(\phi)\\
&\lesssim \left(\frac{\Psi(R)}{\Phi(R)}\int_{B(x_0,8A_{PI}A_{cap}R)}\mathrm{d}\Gamma(f)+\frac{1}{\Phi(R)}\int_{B(x_0,A_{cap}R)}|f|^p\mathrm{d} m\right)\frac{\Phi(R)}{\Psi(R)}\\
&\le\int_{B(x_0,8A_{PI}A_{cap}R)}\mathrm{d}\Gamma(f)+\frac{1}{\Psi(R)}\int_{B(x_0,8A_{PI}A_{cap}R)}|f|^p\mathrm{d} m.
\end{align*}
\end{proof}

\end{document}
